\documentclass[12pt,a4paper]{article}

\makeatletter
\usepackage[utf8]{inputenc}
\usepackage[T2B]{fontenc}
\usepackage[english,russian]{babel}

\usepackage{geometry}
\geometry{
  a4paper,
  left=2.54cm,
  right=2.54cm,
  top=2.54cm,
  bottom=2.54cm,
  headheight=0mm,
  headsep=0mm,
  footnotesep=0mm,
  footskip=0mm,
}

\usepackage{nopageno}

\usepackage{tempora}
\renewcommand{\ttdefault}{cmtt}

\usepackage{titlesec}
\titleformat*{\section}{\large\bfseries}
\titlespacing\section{0pt}{12pt plus 4pt minus 2pt}{0pt plus 2pt minus 2pt}
\setcounter{secnumdepth}{0}

\usepackage{enumitem}

\frenchspacing

\usepackage{amsmath}
\usepackage{amssymb}
\usepackage{hyperref}
\usepackage{longtable}
\usepackage[table]{xcolor}
\usepackage{array}
\usepackage{color}
\usepackage{xcolor}

\usepackage{hyperref}

\newcommand{\MYhref}[3][blue]{\href{#2}{\color{#1}{#3}}}%

\usepackage{listings}
\usepackage{alltt}
\usepackage{csquotes}
\DeclareQuoteStyle{russian}
	{\guillemotleft}{\guillemotright}[0.025em]
	{\quotedblbase}{\textquotedblleft}
\ExecuteQuoteOptions{style=russian}

\usepackage{graphicx}

\lstset{tabsize=2,
	breaklines,
	columns=fullflexible,
	flexiblecolumns,
	numbers=left,
	numberstyle={\footnotesize},
	extendedchars,
	inputencoding=utf8}

\def\@xobeysp{ }
\def\verbatim@processline{\hspace{1.2cm}\raggedright\the\verbatim@line\par}

\parindent=0pt
\parskip=3pt

\definecolor{lightgray}{gray}{0.9}

\renewcommand{\thesubsection}{\arabic{subsection}}

\lstdefinestyle{customc}{
  belowcaptionskip=1\baselineskip,
  breaklines=true,
  frame=L,
  xleftmargin=\parindent,
  language=C,
  showstringspaces=false,
  basicstyle=\footnotesize\ttfamily,
  keywordstyle=\bfseries\color{green!40!black},
  commentstyle=\itshape\color{gray},
  identifierstyle=\color{black},
  stringstyle=\color{blue},
}

\lstdefinestyle{customasm}{
  belowcaptionskip=1\baselineskip,
  frame=L,
  xleftmargin=\parindent,
  language=[x86masm]Assembler,
  basicstyle=\footnotesize\ttfamily,
  commentstyle=\itshape\color{purple!40!black},
}

\lstset{escapechar=@,style=customc}
\makeatother

% Style for terminal output and traces: no C highlighting, no numbers.
\lstdefinestyle{console}{
  language={},
  numbers=none,
  frame=single,
  breaklines=true,
  columns=fullflexible,
  keepspaces=true,
  basicstyle=\footnotesize\ttfamily,
  identifierstyle=\color{black},
}

% listings cannot read multibyte UTF-8: listingsutf8 recodes the source files
% (Cyrillic comments) to KOI8-R on the fly.
\usepackage{listingsutf8}
\lstset{inputencoding=utf8/koi8-r}

\hypersetup{hidelinks}

% Monospaced identifiers cannot be hyphenated; let TeX stretch the spaces
% instead of pushing them into the margin.
\setlength{\emergencystretch}{3em}

\usepackage{caption}
\usepackage{needspace}
\usepackage{float}
\usepackage{multirow}
\captionsetup{font=small}

% Every line number of a source listing becomes a link target "<file>.<line>".
\newcommand{\lstlinetarget}[1]{\raisebox{\baselineskip}[0pt][0pt]{\hypertarget{\lstname.#1}{}}#1}
\lstset{numberstyle=\footnotesize\lstlinetarget}

% Code references: \deffn{key}{file}{line} declares where a definition starts,
% \fn{key} prints the key as a link, \coderef{key}{text} prints custom text.
% The line numbers below must match the sources in ../src.
\newcommand{\deffn}[3]{\expandafter\def\csname cref@#1\endcsname{#2.#3}}
\newcommand{\coderef}[2]{%
  \ifcsname cref@#1\endcsname
    \hyperlink{\csname cref@#1\endcsname}{\textcolor{blue}{\texttt{#2}}}%
  \else
    \errmessage{Unknown code reference: #1}%
  \fi}
\newcommand{\fn}[1]{\coderef{#1}{\detokenize{#1}}}

\deffn{main.c}{main.c}{1}

\deffn{morph.c}{morph.c}{1}

\deffn{morph.h}{morph.h}{1}

\deffn{matrix.c}{matrix.c}{1}

\deffn{matrix.h}{matrix.h}{1}

\deffn{Makefile}{Makefile}{1}


\deffn{kMaxThreads}{main.c}{21}
\deffn{parse_options}{main.c}{92}
\deffn{load_matrix}{main.c}{166}
\deffn{is_same_file}{main.c}{220}
\deffn{filter}{main.c}{251}

\deffn{matrix_t}{matrix.h}{9}
\deffn{matrix_alloc}{matrix.c}{16}
\deffn{matrix_read}{matrix.c}{79}
\deffn{matrix_write}{matrix.c}{122}

\deffn{barrier_t}{morph.c}{11}
\deffn{kChunksPerThread}{morph.c}{40}
\deffn{morph_task_t}{morph.c}{45}
\deffn{barrier_wait}{morph.c}{90}
\deffn{barrier_break}{morph.c}{111}
\deffn{window_min}{morph.c}{118}
\deffn{window_max}{morph.c}{131}
\deffn{filter_rows}{morph.c}{145}
\deffn{filter_both}{morph.c}{171}
\deffn{run_iteration}{morph.c}{178}
\deffn{worker_main}{morph.c}{201}
\deffn{run_workers}{morph.c}{218}
\deffn{morph_apply}{morph.c}{264}

\begin{document}

\begin{titlepage}
\begin{center}
    \includegraphics[width=\textwidth]{mai_header.jpg}

    \vspace{1cm}
    Министерство науки и высшего образования Российской Федерации\\
    Федеральное государственное бюджетное образовательное учреждение высшего образования\\
    \textbf{Московский авиационный институт}\\
    (Национальный исследовательский университет)

    \vspace{1.5cm}
    Институт № 8 \enquote{Компьютерные науки и прикладная математика}\\
    Кафедра 806

    \vspace{2cm}
    {\large\textbf{Отчёт по лабораторной работе №3}}\\[0.3cm]
    {\textbf{Дисциплина}}: \enquote{Операционные системы и архитектура компьютера}\\[0.3cm]
    {\textbf{Тема}}: \enquote{Работа с потоками}\\[0.3cm]
    {\textbf{Вариант}}: 12
\end{center}

\vspace{2cm}
\begin{flushright}
    {\textbf{Выполнил:}}\\
    студент группы М8О-206БВ-25\\
    Ширанков М.С.\\

    {\textbf{Проверил:}}\\
    Ассистент Центра ТОП-ИТ\\
    Ефимов Д.А.
\end{flushright}

\vfill
\begin{center}
    Москва 2026
\end{center}
\end{titlepage}

\tableofcontents
\newpage

\section{Цель работы}
Приобрести практические навыки:
\begin{itemize}
    \item управления потоками в ОС: создания потоков (\texttt{pthread\_create}),
          ожидания их завершения (\texttt{pthread\_join}), ограничения их числа;
    \item обеспечения синхронизации между потоками: мьютексов, условных
          переменных, барьеров и атомарных операций.
\end{itemize}

\section{Задание}
Составить программу на языке Си, обрабатывающую данные в многопоточном режиме.
При обработке использовать стандартные средства создания потоков операционной
системы. Ограничение максимального количества потоков, работающих в один момент
времени, должно задаваться ключом запуска программы. Необходимо уметь
продемонстрировать количество потоков, используемое программой, стандартными
средствами операционной системы. В отчёте привести исследование зависимости
ускорения и эффективности алгоритма от входных данных и количества потоков и
объяснить полученные результаты.

\textbf{Вариант 12:} наложить $K$ раз фильтры эрозии и наращивания на матрицу,
состоящую из вещественных чисел. На выходе получается две результирующие матрицы.

\section{Ход работы}
\subsection{Структура программы}
Программа \texttt{morph} состоит из трёх модулей (полный исходный код приведён
в приложении~А):
\begin{itemize}
    \item \fn{main.c} "--- разбор ключей запуска, ввод и вывод матриц, замер
          времени;
    \item \fn{matrix.c}, \fn{matrix.h} "--- матрица \fn{matrix_t}: выделение памяти,
          чтение и запись в текстовом формате, генерация случайной матрицы;
    \item \fn{morph.c}, \fn{morph.h} "--- многопоточное применение фильтров:
          потоки, барьер, распределение работы.
\end{itemize}

Сборка выполняется по \fn{Makefile} из корня лабораторной работы, замеры для
раздела исследования "--- скриптом \texttt{bench/bench.py}:
\begin{lstlisting}[style=console]
make
./morph -t 4 -k 3 -w 5 -i input.txt -e eroded.txt -d dilated.txt
python3 bench/bench.py
\end{lstlisting}

\needspace{18\baselineskip}
\subsection{Ключи запуска}
Формат матрицы: два целых числа "--- количество строк и столбцов, затем все
элементы по строкам через любые пробельные символы (пробелы, табуляции, перевод строки и т.д.). В этом же формате
выводятся результаты.

\begin{longtable}{|>{\raggedright\arraybackslash}p{0.19\textwidth}|>{\raggedright\arraybackslash}p{0.71\textwidth}|}
\hline
\rowcolor{lightgray}
\textbf{Ключ} & \textbf{Назначение} \\
\hline
\endhead
\texttt{-t THREADS} & \textbf{обязательный}: максимальное число потоков, работающих
                      одновременно, включая главный (1--\fn{kMaxThreads}) \\
\hline
\texttt{-k K} & сколько раз применить каждый фильтр (по умолчанию 1) \\
\hline
\texttt{-w WINDOW} & нечётная сторона квадратного окна (по умолчанию 3) \\
\hline
\texttt{-i FILE} & файл с матрицей; без ключа матрица читается из \texttt{stdin} \\
\hline
\texttt{-g RxC} & сгенерировать случайную матрицу $R \times C$ вместо чтения \\
\hline
\texttt{-s SEED} & зерно генератора для \texttt{-g} (по умолчанию 1) \\
\hline
\texttt{-e FILE},\newline\texttt{-d FILE} & файлы для результата эрозии и наращивания;
                       без ключа результат выводится в \texttt{stdout} \\
\hline
\texttt{-S} & статическое распределение строк между потоками
              (по умолчанию "--- динамическое, см. ниже) \\
\hline
\texttt{-q} & не выводить результаты, для которых не задан файл \\
\hline
\texttt{-v} & вывести в \texttt{stderr} PID, число потоков и время фильтрации \\*
\hline
\end{longtable}

\subsection{Алгоритм}
Эрозия и наращивание "--- операции математической морфологии. Для матрицы $A$
размера $N \times M$ и квадратного окна стороны $w = 2r + 1$:
\[
    (A \ominus W)_{ij} = \min_{\substack{|a| \le r,\ |b| \le r \\ 0 \le i+a < N,\ 0 \le j+b < M}} A_{i+a,\,j+b},
    \qquad
    (A \oplus W)_{ij} = \max_{\substack{|a| \le r,\ |b| \le r \\ 0 \le i+a < N,\ 0 \le j+b < M}} A_{i+a,\,j+b}.
\]
У края матрицы окно обрезается (\fn{filter_rows}): учитываются только
существующие элементы. Это равносильно дополнению матрицы значением $+\infty$
для эрозии и $-\infty$ для наращивания, поэтому края не искажаются.

Фильтры применяются $K$ раз: $E_0 = D_0 = A$, $E_k = E_{k-1} \ominus W$,
$D_k = D_{k-1} \oplus W$; результат "--- $E_K$ и $D_K$. Минимум и максимум по окну
считаются по определению (\fn{window_min}, \fn{window_max}), поэтому трудоёмкость
равна $\Theta(K \cdot N \cdot M \cdot w^2)$ сравнений на каждый фильтр, а размер окна
напрямую управляет объёмом вычислений на один элемент.

\subsection{Распараллеливание}
Внутри одной итерации каждый элемент результата зависит только от матрицы
предыдущей итерации, поэтому все элементы можно считать независимо. Матрица
делится между потоками по строкам; каждую строку результата записывает ровно
один поток, и для записи блокировки не нужны. Обе цепочки, эрозия и наращивание,
обрабатываются одними и теми же потоками: поток, взявший строки, считает в них
оба фильтра.

\textbf{Двойная буферизация.} Для каждого фильтра есть два буфера: итерация $i$
читает буфер итерации $i-1$ и пишет в \texttt{out[i \% 2]} (\fn{morph_task_t},
\fn{filter_both}).
Первая итерация читает исходную матрицу напрямую, а буферы назначаются так, чтобы
последняя, $K$-я, запись попала сразу в выходную матрицу (\fn{morph_apply}).
Копирования между итерациями нет.

\textbf{Барьер.} Итерации зависят друг от друга, и опасностей две:
\begin{itemize}
    \item \enquote{чтение до записи}: строки на границе участка читают строки
          соседнего участка, и итерация $i+1$ не может начаться, пока сосед не
          дописал их на итерации $i$;
    \item \enquote{запись до окончания чтения}: итерация $i+1$ пишет в буфер
          \texttt{out[(i+1) \% 2]}, который на итерации $i$ ещё читают другие
          потоки как результат итерации $i-1$.
\end{itemize}
Обе снимает один барьер между итерациями: никакой поток не начинает итерацию
$i+1$, пока все не закончили итерацию $i$. После последней итерации барьер не
нужен "--- его роль играет \texttt{pthread\_join}.

В macOS нет \texttt{pthread\_barrier\_t}, поэтому барьер \fn{barrier_t} реализован
на мьютексе и условной переменной (\fn{barrier_wait}). Поток увеличивает счётчик
пришедших; последний обнуляет его, увеличивает номер поколения
(\texttt{generation}) и будит остальных через \texttt{pthread\_cond\_broadcast}.
Остальные ждут в цикле, пока не сменится поколение. Номер поколения нужен по двум
причинам: \texttt{pthread\_cond\_wait} допускает ложные пробуждения, а поток,
быстро прошедший барьер, может прийти к нему снова раньше, чем проснулись
остальные, "--- проверка счётчика вместо поколения в этом случае сломалась бы.

\textbf{Ограничение числа потоков.} Ключ \texttt{-t} задаёт число $T$; создаётся
$\min(T, N)$ вычисляющих потоков, так как делить работу мельче строки нет смысла.
Главный поток не простаивает в \texttt{pthread\_join}, а работает как поток номер 0
(\fn{run_workers}), поэтому в процессе ровно $\min(T, N)$ потоков. Потоки создаются
один раз и живут все $K$ итераций; создавать их заново на каждой итерации было бы
дороже.

Перед вычислениями все потоки ждут на стартовом барьере (\fn{worker_main}). Если
\texttt{pthread\_create} не удался, уже запущенные потоки ждали бы на барьере
вечно: число участников не набралось бы. Поэтому барьер можно
\enquote{сломать} (\fn{barrier_break}): ожидающие просыпаются, получают $-1$ и
завершаются, а главный поток дожидается их и сообщает об ошибке.

\textbf{Распределение строк.} Реализованы две стратегии:
\begin{itemize}
    \item \textbf{статическая} (ключ \texttt{-S}): поток $t$ из $T$ получает
          полосу строк $[\lfloor tN/T \rfloor, \lfloor (t+1)N/T \rfloor)$; полосы
          отличаются по высоте не больше чем на одну строку;
    \item \textbf{динамическая} (по умолчанию): матрица делится на блоки по
          $\lceil N / (16 T) \rceil$ строк (\fn{kChunksPerThread}), а поток, закончив
          блок, берёт номер следующего атомарной операцией
          \texttt{atomic\_fetch\_add} над общим счётчиком. Быстрый поток успевает
          обработать больше блоков, чем медленный.
\end{itemize}
Счётчик блоков после итерации нужно обнулить, но сделать это сразу нельзя: другие
потоки ещё могут обращаться к нему. Поэтому счётчиков два, как и буферов:
итерация $i$ раздаёт блоки из \texttt{next\_chunk[i \% 2]}, а поток 0 в её начале
обнуляет \texttt{next\_chunk[(i+1) \% 2]} (\fn{run_iteration}). Этот счётчик использовался на
итерации $i-1$, и все потоки уже прошли барьер после неё; обнуление станет
видно остальным через барьер после итерации $i$. Порядок обращений к
счётчику обеспечивает барьер на мьютексе, поэтому для самой атомарной операции
достаточно \texttt{memory\_order\_relaxed}.

\needspace{10\baselineskip}
\section{Используемые системные вызовы и функции}
\begin{longtable}{|>{\raggedright\arraybackslash}p{0.27\textwidth}|>{\raggedright\arraybackslash}p{0.42\textwidth}|>{\raggedright\arraybackslash}p{0.21\textwidth}|}
\hline
\rowcolor{lightgray}
\textbf{Вызов} & \textbf{Назначение в программе} & \textbf{Где в коде} \\
\hline
\endhead
\texttt{pthread\_create} & создание вычисляющих потоков & \fn{run_workers} \\
\hline
\texttt{pthread\_join} & ожидание завершения потоков и освобождение их ресурсов &
                         \fn{run_workers} \\
\hline
\texttt{pthread\_mutex\_init}, \texttt{\_lock}, \texttt{\_unlock}, \texttt{\_destroy} &
    защита состояния барьера & \fn{barrier_t}, \fn{barrier_wait} \\
\hline
\texttt{pthread\_cond\_init}, \texttt{\_wait}, \texttt{\_broadcast}, \texttt{\_destroy} &
    ожидание остальных потоков на барьере и их пробуждение &
    \fn{barrier_wait}, \fn{barrier_break} \\
\hline
\texttt{atomic\_fetch\_add}, \texttt{atomic\_store} & выдача блоков строк при
    динамическом распределении & \fn{run_iteration} \\
\hline
\texttt{clock\_gettime} & замер времени фильтрации по монотонным часам
                          (\texttt{CLOCK\_MONOTONIC}) & \fn{filter} \\
\hline
\texttt{malloc}, \texttt{free} & память под матрицы и промежуточные буферы &
                                 \fn{matrix_alloc}, \fn{morph_apply} \\
\hline
\texttt{fopen}, \texttt{fscanf}, \texttt{fprintf}, \texttt{fclose} & чтение
    исходной и запись результирующих матриц & \fn{load_matrix},
    \fn{matrix_read}, \fn{matrix_write} \\
\hline
\texttt{fstat} & проверка, что результаты не записываются в один файл &
                 \fn{is_same_file} \\
\hline
\texttt{getopt} & разбор ключей запуска & \fn{parse_options} \\*
\hline
\end{longtable}

На уровне ядра \texttt{pthread\_create} в Linux выполняется системным вызовом
\texttt{clone3} с флагами \texttt{CLONE\_VM}, \texttt{CLONE\_THREAD} и др.:
новый поток разделяет с создателем адресное пространство, дескрипторы и
обработчики сигналов. Ожидание на мьютексе и условной переменной при
конфликте превращается в вызов \texttt{futex}; без конфликта мьютекс
захватывается атомарной операцией в пространстве пользователя, без системного
вызова.

\section{Демонстрация работы}
\subsection{Пример фильтрации}
Матрица $4 \times 5$, окно $3 \times 3$. При $K = 1$ оба результата
выводятся в \texttt{stdout} (сначала эрозия, затем наращивание), при $K = 2$
"--- в файлы:

\begin{lstlisting}[style=console]
$ cat input.txt
4 5
1 9 2 8 3
7 0.5 6 4 5
2 3 -1 9 6
8 4 7 1 2
$ ./morph -t 2 -w 3 -k 1 -i input.txt
4 5
0.5 0.5 0.5 2 3
0.5 -1 -1 -1 3
0.5 -1 -1 -1 1
2 -1 -1 -1 1
4 5
9 9 9 8 8
9 9 9 9 9
8 8 9 9 9
8 8 9 9 9
$ ./morph -t 2 -w 3 -k 2 -i input.txt -e eroded.txt -d dilated.txt -v
pid 71752: filtering 4x5...
matrix=4x5 window=3 iterations=2 threads=2 split=dynamic time_ms=0.084
$ cat eroded.txt dilated.txt
4 5
-1 -1 -1 -1 -1
-1 -1 -1 -1 -1
-1 -1 -1 -1 -1
-1 -1 -1 -1 -1
4 5
9 9 9 9 9
9 9 9 9 9
9 9 9 9 9
8 9 9 9 9
\end{lstlisting}

Например, элемент $(0, 0)$ после эрозии равен $\min(1, 9, 7, 0{,}5) = 0{,}5$:
окно в углу обрезано до $2 \times 2$. Эрозия \enquote{растягивает} минимумы, а
наращивание "--- максимумы: каждый проход сдвигает их границу на $r = 1$
элемент. От минимума $-1$ в позиции $(2, 2)$ до любого элемента матрицы не
больше двух шагов по строкам и столбцам, поэтому за два прохода он заполнил
всю матрицу. А до элемента $(3, 0)$ от обеих девяток, $(0, 1)$ и $(2, 3)$,
три шага, поэтому после наращивания там осталось 8.

\subsection{Количество потоков средствами ОС}
Программа разрабатывалась на macOS (Apple M1). Запускаем долгую фильтрацию с
ограничением в 6 потоков и смотрим на процесс из другого терминала. Утилита
\texttt{ps} с ключом \texttt{-M} выводит по строке на каждый поток, а
\texttt{top} "--- число потоков в столбце \texttt{\#TH}:
\begin{lstlisting}[style=console]
$ ./morph -t 6 -g 3000x3000 -w 9 -k 30 -q -v &
pid 74665: filtering 3000x3000...
$ ps -M -p 74665
USER       PID   TT   %CPU STAT PRI     STIME     UTIME COMMAND
eblackbu 74665   ??   66.3 R    31T   0:00.04   0:00.79 ./morph -t 6 -g 3000x3000 -w 9 -k 30 -q -v
         74665        70.0 R    31T   0:00.04   0:00.80
         74665        67.9 R    31T   0:00.03   0:00.77
         74665        66.8 R    31T   0:00.03   0:00.77
         74665        65.7 R    31T   0:00.03   0:00.75
         74665        67.5 R    31T   0:00.03   0:00.78
$ top -l 1 -pid 74665 -stats pid,command,cpu,threads
PID    COMMAND %CPU #TH
74665  morph   0.0  6/6
\end{lstlisting}
Потоков ровно шесть: главный и пять созданных, все в состоянии \texttt{R}
(выполняется) с почти одинаковым процессорным временем \texttt{UTIME} "---
работа распределена равномерно. (\texttt{top -l 1} делает один снимок, поэтому
загрузка в нём не измеряется и показана как 0.)

То же в Linux (контейнер Debian, ядро 6.12, \texttt{gcc}~14.2). Здесь каждый
поток "--- отдельная задача ядра со своим идентификатором \texttt{SPID}
(\texttt{TID}):
\begin{lstlisting}[style=console]
$ ps -T -p 611
  PID  SPID TTY          TIME CMD
  611   611 ?        00:00:01 morph
  611   613 ?        00:00:01 morph
  611   614 ?        00:00:01 morph
  611   615 ?        00:00:01 morph
  611   616 ?        00:00:01 morph
  611   617 ?        00:00:01 morph
$ ps -o pid,nlwp,cmd -p 611
  PID NLWP CMD
  611    6 ./morph -t 6 -g 3000x3000 -w 9 -k 30 -q -v
$ grep Threads /proc/611/status
Threads:	6
$ ls /proc/611/task
611  613  614  615  616  617
\end{lstlisting}
Также число потоков видно в \texttt{top -H} и \texttt{htop}. Если строк в
матрице меньше, чем задано \texttt{-t}, потоков создаётся по числу строк.

\subsection{Трассировка системных вызовов}
\begin{lstlisting}[style=console]
strace -f -e trace=clone3,futex,exit,exit_group ./morph -t 3 -g 600x600 -k 2 -q
\end{lstlisting}
Трасса (PID 625 "--- главный поток, 626 и 627 "--- созданные; длинные
аргументы сокращены):
\begin{lstlisting}[style=console]
625   clone3({flags=CLONE_VM|CLONE_FS|CLONE_FILES|CLONE_SIGHAND|CLONE_THREAD|
              CLONE_SYSVSEM|CLONE_SETTLS|CLONE_PARENT_SETTID|CLONE_CHILD_CLEARTID,
              ..., stack_size=0x80ea60, ...} => {parent_tid=[626]}, 88) = 626
626   futex(0xffffc4d83010, FUTEX_WAIT_BITSET_PRIVATE|..., 0, ...) <unfinished ...>
625   clone3({flags=...|CLONE_THREAD|..., ...} => {parent_tid=[627]}, 88) = 627
625   futex(0xffffc4d83010, FUTEX_WAIT_BITSET_PRIVATE|..., 0, ...) <unfinished ...>
627   futex(0xffffc4d83010, FUTEX_WAKE_PRIVATE, 2147483647) = 2
625   <... futex resumed>)              = 0
626   <... futex resumed>)              = 0
625   futex(0xffffc4d82fc0, FUTEX_WAKE_PRIVATE, 1 <unfinished ...>
626   futex(0xffffc4d82fc0, FUTEX_WAKE_PRIVATE, 1 <unfinished ...>
...
627   futex(0xffffc4d83014, FUTEX_WAIT_BITSET_PRIVATE|..., 0, ...) <unfinished ...>
626   futex(0xffffc4d83014, FUTEX_WAIT_BITSET_PRIVATE|..., 0, ...) <unfinished ...>
625   futex(0xffffc4d83014, FUTEX_WAKE_PRIVATE, 2147483647) = 2
...
625   futex(0xffff84b9f250, FUTEX_WAIT_BITSET|..., 626, ...) <unfinished ...>
627   exit(0)                           = ?
626   exit(0 <unfinished ...>
625   <... futex resumed>)              = 0
625   exit_group(0)                     = ?
\end{lstlisting}

По трассе видно:
\begin{itemize}
    \item \texttt{pthread\_create} "--- это \texttt{clone3} с флагами
          \texttt{CLONE\_VM} (общая память), \texttt{CLONE\_FILES} (общие
          дескрипторы), \texttt{CLONE\_SIGHAND} и \texttt{CLONE\_THREAD} (тот же
          процесс), под стек нового потока выделено около 8~МБ
          (\texttt{stack\_size=0x80ea60});
    \item барьер проходится дважды: стартовый и между двумя итерациями. На
          первом последним пришёл поток 627, на втором "--- главный поток 625.
          Пришедшие раньше спят в \texttt{FUTEX\_WAIT} на условной переменной
          (адреса \texttt{...3010} и \texttt{...3014} "--- её внутренний счётчик),
          последний будит всех одним \texttt{FUTEX\_WAKE} с числом
          $2^{31}-1$ (\texttt{pthread\_cond\_broadcast}), ядро сообщает, что
          проснулись два потока;
    \item проснувшиеся потоки одновременно захватывают мьютекс барьера
          (\texttt{...2fc0}), отсюда короткие \texttt{FUTEX\_WAKE} на нём.
          Без конфликта мьютекс системных вызовов не делает "--- в трассе нет
          захватов при чтении счётчика блоков, он атомарный;
    \item \texttt{pthread\_join} ждёт поток на \texttt{futex} по адресу его
          идентификатора: ядро обнуляет его и будит ожидающего при выходе потока
          (флаг \texttt{CLONE\_CHILD\_CLEARTID});
    \item созданные потоки завершаются вызовом \texttt{exit}, который
          завершает только поток, а главный "--- \texttt{exit\_group}, завершающим
          весь процесс.
\end{itemize}

Сводная статистика \texttt{strace -f -c} для $1000 \times 1000$, $K = 10$, 4
потока: 93 вызова \texttt{futex} (99\% системного времени, это время ожидания
на барьере) и 3 вызова \texttt{clone3}. Вызовов на элемент матрицы нет совсем:
вся фильтрация выполняется в пространстве пользователя, ядро участвует только в
создании потоков и в их ожидании.


\section{Исследование ускорения и эффективности}
% Earlier longtables have no captions but still advance the counter.
\setcounter{table}{0}
Ускорение и эффективность на $p$ потоках:
\[
    S_p = \frac{T_1}{T_p}, \qquad E_p = \frac{S_p}{p},
\]
где $T_p$ "--- время фильтрации на $p$ потоках (включая создание и завершение
потоков), $T_1$ "--- время на одном потоке.

Замеры выполнены на MacBook Air (Apple M1: 4 производительных ядра и
4 энергоэффективных, примерно в 2--3 раза более медленных). Скрипт
\texttt{bench/bench.py} запускает каждую конфигурацию 9 раз в случайном
порядке и берёт медиану времени; результаты сохраняются в
\texttt{bench/results/*.csv}. Сводная таблица по этим результатам:

\begin{table}[H]
\centering
\caption{Время на одном потоке, ускорение и эффективность}
\label{tab:speedup}
\begin{tabular}{|c|c|r|c|c|c|c||c|c|}
\hline
\rowcolor{lightgray}
\textbf{Параметры} & \textbf{Распр.} & \textbf{$T_1$, мс} & \textbf{$S_2$} & \textbf{$S_4$} &
\textbf{$S_8$} & \textbf{$S_{16}$} & \textbf{$E_4$} & \textbf{$E_8$} \\
\hline
\multicolumn{9}{|l|}{\textit{Размер матрицы $N \times N$, окно $3 \times 3$, $K = 10$}} \\
\hline
\multirow{2}{*}{$N = 50$} & стат. & 0{,}52 & 1{,}34 & 1{,}66 & 0{,}94 & 0{,}57 & 0{,}41 & 0{,}12 \\
 & дин. & 0{,}52 & 1{,}49 & 1{,}59 & 0{,}78 & 0{,}49 & 0{,}40 & 0{,}10 \\
\hline
\multirow{2}{*}{$N = 100$} & стат. & 2{,}09 & 1{,}77 & 2{,}57 & 2{,}18 & 1{,}91 & 0{,}64 & 0{,}27 \\
 & дин. & 2{,}09 & 1{,}78 & 2{,}91 & 2{,}19 & 1{,}65 & 0{,}73 & 0{,}27 \\
\hline
\multirow{2}{*}{$N = 200$} & стат. & 8{,}35 & 1{,}93 & 3{,}46 & 4{,}05 & 3{,}50 & 0{,}86 & 0{,}51 \\
 & дин. & 8{,}35 & 1{,}94 & 3{,}55 & 4{,}14 & 3{,}33 & 0{,}89 & 0{,}52 \\
\hline
\multirow{2}{*}{$N = 500$} & стат. & 52{,}1 & 1{,}98 & 3{,}86 & 4{,}79 & 4{,}90 & 0{,}96 & 0{,}60 \\
 & дин. & 52{,}2 & 1{,}94 & 3{,}84 & 5{,}62 & 5{,}43 & 0{,}96 & 0{,}70 \\
\hline
\multirow{2}{*}{$N = 1000$} & стат. & 210 & 1{,}99 & 3{,}86 & 5{,}26 & 5{,}11 & 0{,}96 & 0{,}66 \\
 & дин. & 210 & 1{,}96 & 3{,}85 & 5{,}74 & 5{,}91 & 0{,}96 & 0{,}72 \\
\hline
\multirow{2}{*}{$N = 2000$} & стат. & 846 & 1{,}99 & 3{,}81 & 4{,}64 & 4{,}83 & 0{,}95 & 0{,}58 \\
 & дин. & 845 & 1{,}95 & 3{,}76 & 5{,}40 & 5{,}90 & 0{,}94 & 0{,}68 \\
\hline
\multicolumn{9}{|l|}{\textit{Размер окна $w$, матрица $1000 \times 1000$, $K = 4$}} \\
\hline
\multirow{2}{*}{$w = 1$} & стат. & 20{,}2 & 1{,}85 & 3{,}28 & 2{,}81 & 2{,}81 & 0{,}82 & 0{,}35 \\
 & дин. & 20{,}2 & 1{,}86 & 3{,}29 & 2{,}87 & 2{,}57 & 0{,}82 & 0{,}36 \\
\hline
\multirow{2}{*}{$w = 3$} & стат. & 87{,}1 & 2{,}00 & 3{,}82 & 5{,}00 & 5{,}01 & 0{,}96 & 0{,}63 \\
 & дин. & 86{,}6 & 1{,}96 & 3{,}70 & 5{,}45 & 5{,}30 & 0{,}93 & 0{,}68 \\
\hline
\multirow{2}{*}{$w = 5$} & стат. & 190 & 2{,}00 & 3{,}93 & 5{,}22 & 5{,}25 & 0{,}98 & 0{,}65 \\
 & дин. & 193 & 1{,}97 & 3{,}73 & 6{,}02 & 6{,}06 & 0{,}93 & 0{,}75 \\
\hline
\multirow{2}{*}{$w = 9$} & стат. & 816 & 2{,}04 & 4{,}04 & 6{,}01 & 6{,}04 & 1{,}01 & 0{,}75 \\
 & дин. & 816 & 2{,}00 & 3{,}99 & 6{,}33 & 6{,}67 & 1{,}00 & 0{,}79 \\
\hline
\multirow{2}{*}{$w = 15$} & стат. & 3150 & 2{,}00 & 3{,}98 & 6{,}94 & 6{,}82 & 1{,}00 & 0{,}87 \\
 & дин. & 3153 & 1{,}97 & 3{,}91 & 7{,}29 & 7{,}50 & 0{,}98 & 0{,}91 \\
\hline
\end{tabular}
\end{table}

Выводы:
\begin{itemize}
    \item \textbf{До 4 потоков} ускорение почти линейное ($E_4 \approx 0{,}94$--$1{,}0$
          для $N \ge 500$ и $w \ge 3$): потоки пишут в разные строки, а
          синхронизация "--- один барьер на итерацию.
    \item \textbf{Малые матрицы} ускоряются плохо: создание потока стоит
          около 14~мкс, проход барьера "--- десятки микросекунд, а итерация для
          $50 \times 50$ длится всего около 52~мкс. Поэтому уже на 8 потоках
          программа медленнее последовательной ($S_8 < 1$).
    \item \textbf{От 5 до 8 потоков} подключаются медленные E-ядра. При
          статическом распределении время итерации определяет поток на
          E-ядре, а остальные ждут его на барьере; динамическое распределение
          отдаёт быстрым потокам больше блоков и выигрывает
          (для $N = 1000$: $S_8 = 5{,}74$ против 5{,}26).
    \item \textbf{Больше 8 потоков} "--- больше, чем ядер: ускорение не растёт,
          а эффективность падает.
    \item \textbf{Чем больше окно}, тем больше вычислений приходится на одну
          синхронизацию и тем ближе ускорение к идеальному: при $w = 15$
          получено $S_8 = 7{,}29$. Окно $1 \times 1$ "--- это просто копирование, оно
          упирается в память, и ускорение не превышает 3{,}3.
\end{itemize}

\section{Итог работы}
В ходе лабораторной работы реализована программа многократного применения
фильтров эрозии и наращивания к матрице вещественных чисел с использованием
потоков POSIX. Число одновременно работающих потоков задаётся ключом
\texttt{-t}, а их фактическое количество продемонстрировано средствами macOS
(\texttt{ps -M}, \texttt{top}) и Linux (\texttt{ps -T}, \texttt{/proc}). Получены
практические навыки создания и ожидания потоков (\texttt{pthread\_create},
\texttt{pthread\_join}) и их синхронизации: собственный барьер на мьютексе и
условной переменной, атомарный счётчик для распределения работы.

\newpage
\phantomsection
\addcontentsline{toc}{section}{Приложение А. Исходный код программы}
\begin{center}
    {\large\bfseries ПРИЛОЖЕНИЕ А}\\[0.2cm]
    {\bfseries Исходный код программы}
\end{center}

Исходный код также доступен в
\MYhref{https://github.com/evdakCherpak/os_mai_labs}{репозитории на GitHub}.

\phantomsection
\addcontentsline{toc}{subsection}{А.1. main.c}
\subsection*{А.1. main.c}
\lstinputlisting[name=main.c]{../src/main.c}

\phantomsection
\addcontentsline{toc}{subsection}{А.2. morph.h}
\subsection*{А.2. morph.h}
\lstinputlisting[name=morph.h]{../src/morph.h}

\phantomsection
\addcontentsline{toc}{subsection}{А.3. morph.c}
\subsection*{А.3. morph.c}
\lstinputlisting[name=morph.c]{../src/morph.c}

\phantomsection
\addcontentsline{toc}{subsection}{А.4. matrix.h}
\subsection*{А.4. matrix.h}
\lstinputlisting[name=matrix.h]{../src/matrix.h}

\phantomsection
\addcontentsline{toc}{subsection}{А.5. matrix.c}
\subsection*{А.5. matrix.c}
\lstinputlisting[name=matrix.c]{../src/matrix.c}

\phantomsection
\addcontentsline{toc}{subsection}{А.6. Makefile}
\subsection*{А.6. Makefile}
% escapechar=@ from log-style would break on $@ in the Makefile
\lstinputlisting[name=Makefile,language=make,escapechar={}]{../Makefile}

\end{document}
